% !TeX program = pdflatex
% Related-work section for the Cascade Engine paper.
% Self-contained (compiles standalone) but written so the \section can be
% \input{} directly into the main manuscript.

\documentclass[11pt]{article}
\usepackage[margin=1in]{geometry}
\usepackage{amsmath, amssymb}
\usepackage{hyperref}
\usepackage{microtype}
\title{Related Work}
\author{}
\date{}
\begin{document}
\maketitle

\section{Related Work}
\label{sec:related}

We organize prior work into three threads --- \emph{routing}, \emph{cascading},
and their recent \emph{unification} --- followed by the non-stationary
online-learning literature our analysis builds on. We state our delta against
the closest work explicitly rather than leaving it implicit.

\paragraph{Routing: one model per query.}
Routing methods train a predictor that maps a query to a single model, trading
quality against cost without ever calling more than one model. RouteLLM
(Ong et al.\ 2024) learns a binary strong/weak router from Chatbot Arena
preference data using matrix factorization, a BERT classifier, or a causal-LLM
classifier, and is the de-facto open baseline; we reproduce its released BERT
checkpoint rather than a proxy. Hybrid LLM (Ding et al.\ 2024) trains a router
on a quality gap signal and exposes a cost knob. ZOOTER (Lu et al.\ 2023)
distills reward-model supervision into a lightweight router. RouterBench
(Hu et al.\ 2024) standardizes evaluation across routers, and FORC
(\v{S}akota et al.\ 2024) frames routing as a meta-model assignment problem.
All of these are \emph{static}: the router is trained offline and frozen, and
the analysis (where present) is about the learned predictor's accuracy, not
about regret under a changing environment.

\paragraph{Cascading: sequential escalation.}
Cascading runs models from cheapest to most capable, stopping when an answer is
deemed good enough. FrugalGPT (Chen et al.\ 2023) introduced the LLM-cascade:
a learned scorer judges each response and triggers escalation; we reproduce
this with a public reward model as the scorer rather than a hand-tuned
heuristic. AutoMix (Madaan et al.\ 2023) uses a self-verification signal and a
meta-verifier to decide escalation, and FrugalGPT-style deferral has been
analyzed under learning-to-defer framing. These methods are effective but,
prior to 2025, lacked optimality proofs and could not be combined with routing
in a principled way.

\paragraph{Unified routing and cascading.}
The closest work to ours is Dekoninck, Baader, and Vechev (2025), \emph{``A
Unified Approach to Routing and Cascading for LLMs''} (ICML 2025). They derive
a \emph{per-query optimal} cascading strategy, prove optimality of an existing
routing strategy, and unify both into \emph{cascade routing}, identifying
quality-estimator quality as the decisive factor. Their result is
decision-theoretic and \emph{static}: \emph{given} a quality estimator (with a
known uncertainty covariance) and a cost estimator, they characterize the
optimal one-shot model-selection rule, with hyperparameters fit offline by
Bayesian search over pre-computed responses.

\paragraph{Our delta.}
We do not contest, and we adopt as a baseline, the static optimality of
Dekoninck et al.\ (2025): \emph{their per-query rule is the right target when a
trustworthy quality estimator is available and the environment is stationary.}
Our contribution is orthogonal and addresses precisely the question their
framework assumes away --- their own analysis names good quality estimators as
the critical factor but takes them as given. We study the regime in which the
per-engine quality/reliability signal is (i) \emph{unknown a priori} and (ii)
\emph{non-stationary}, drifting with load, time-of-day, and model deprecation.
In this regime we give Contextual Discounted Thompson Sampling (CD-TS), an
online algorithm with a non-stationary pseudo-regret bound
$\widetilde{O}\big((MK)^{1/3} V_T^{1/3} T^{2/3}\big)$ against the \emph{dynamic}
optimum and a matching lower bound. Informally: Dekoninck et al.\ characterize
\emph{what} the optimal decision is given an estimator; we characterize \emph{how
fast} an agent provably converges to good decisions while \emph{learning} that
estimator online under drift. The two results compose --- their rule is the
per-step oracle our regret is measured against --- and neither subsumes the
other.

\paragraph{Non-stationary bandits and RL.}
Our analysis instantiates the discounted/sliding-window non-stationary bandit
program. Garivier and Moulines (2008) established the $T^{2/3}$ regime for
switching bandits via discounted/windowed UCB; Besbes, Gur, and Zeevi (2014)
gave the variation-budget formulation and the matching
$\Omega(V_T^{1/3}T^{2/3})$ lower bound we cite for tightness. Russac, Vernade,
and Capp\'e (2019) and Trov\`o, Restelli, and Gatti (2020) extended weighted
and sliding-window estimators to linear and Thompson settings respectively;
Cheung, Simchi-Levi, and Zhu (2019) removed the need to know $V_T$ via a
bandit-over-bandit meta-algorithm. For the MDP variant we reference the
minimax episodic-RL bounds of Azar, Osband, and Munos (2017), the model-free
\textsc{UCB-Q} analysis of Jin et al.\ (2018), and the lower bound of
Domingues et al.\ (2021). Our technical novelty is not a new bandit algorithm
but the \emph{instantiation} of this machinery for LLM cascade routing: the
contextual structure is the prompt-complexity binning, the arms are engine
tiers, the reward is the cost--quality composite scored by an external reward
model, and the non-stationarity is the empirically-observed drift in cloud-LLM
reliability.

\paragraph{Benchmarking methodology.}
A secondary contribution is methodological. Prior routing/cascading evaluations
frequently (i) compare against weakened reimplementations of baselines,
(ii) report a single ad-hoc cost-quality scalar rather than a Pareto frontier,
(iii) use the asymptotic normal CI at small seed counts, and (iv) bundle
caching/PII/gatekeeper layers into one method's pipeline but not the baselines'.
We use released baseline checkpoints, a reward-model quality axis independent of
the router's self-report, a true Pareto-frontier analysis, $t$-distribution
confidence intervals at the actual seed count, and an orchestration wrapper that
applies system-level layers uniformly across all routers or to none. Every run
emits a manifest pinning code (git SHA), inputs (dataset SHA-256), toolchain
(package versions), and randomness (seed list) for exact reproduction.

\end{document}
